%% Wave-17 Opus Attack-Heal 11: SC^{ch,top} Heptagon
%% Elite voice: KAPRANOV (higher structure) + ETINGOF (clarity) + COSTELLO (factorisation)
%% Target: /Users/raeez/chiral-bar-cobar-vol2/chapters/theory/sc_chtop_heptagon.tex (2749 lines).
%% Date: 2026-04-24

\providecommand{\SCchtop}{\mathsf{SC}^{\mathrm{ch},\mathrm{top}}}
\providecommand{\FM}{\mathrm{FM}}
\providecommand{\Conf}{\mathrm{Conf}}
\providecommand{\ChirHoch}{\mathrm{ChirHoch}}
\providecommand{\Obs}{\mathsf{Obs}}
\providecommand{\Rep}{\mathrm{Rep}}
\providecommand{\cZ}{\mathcal{Z}}
\providecommand{\cA}{\mathcal{A}}
\providecommand{\cM}{\mathcal{M}}
\providecommand{\cY}{\mathcal{Y}}
\providecommand{\Ran}{\mathrm{Ran}}
\providecommand{\Dmod}{D\text{-}\mathsf{mod}}
\providecommand{\ChirAlg}{\mathsf{ChirAlg}}
\providecommand{\End}{\mathrm{End}}
\providecommand{\Lie}{\mathsf{Lie}}
\providecommand{\Ass}{\mathsf{Ass}}
\providecommand{\id}{\mathrm{id}}
\providecommand{\Hom}{\mathrm{Hom}}
\providecommand{\Dhol}{\mathrm{HolFA}}
\providecommand{\MC}{\mathrm{MC}}
\providecommand{\Eone}{E_1}
\providecommand{\fg}{\mathfrak{g}}
\providecommand{\fsl}{\mathfrak{sl}}

\section*{Wave-17 opus attack-heal 11: seven presentations of \texorpdfstring{$\SCchtop$}{SC\textasciicircum ch,top}}

Reading of \texttt{chiral-bar-cobar-vol2/chapters/theory/sc\_chtop\_heptagon.tex} (2749 lines, Prologue through ``class-$\mathrm M$ heptagon face at $c_{2d} = -214$'') and of the attached Vol~II \texttt{chiral\_higher\_deligne.tex} section~\ref*{sec:chd-braces}. Adversarial programme: six attack axes from the task prompt, five full attack-heal cycles, each ending with an explicit heal; V_k($\fsl_2$) at $k = 1$ used as the simplest nontrivial concrete witness. Writing in the Kapranov-Etingof voice: higher structure first, clarity second; factorisation-algebra pedigree (Costello) retained throughout.

%%%%%%%%%%%%%%%%%%%%%%%%%%%%%%%%%%%%%%%%%%%%%%%%%%%%%%%%%%%%%%%%%%%%%%%%%%%%%
\section{Cycle 0: the state of the heptagon as inscribed (reading log)}
%%%%%%%%%%%%%%%%%%%%%%%%%%%%%%%%%%%%%%%%%%%%%%%%%%%%%%%%%%%%%%%%%%%%%%%%%%%%%

The source chapter presents a specified object --- the two-coloured coloured operad $\SCchtop$ of Definition~\ref*{def:factorization-swiss-cheese-operad} with colours $\{\mathsf{cl}, \mathsf{op}\}$, closed component $\FM_k(\mathbb H)$ fibered over $\Conf_m(\R)$, and empty closed-output component whenever $m \ge 1$ (directional restriction: no open-to-closed operation) --- and exhibits seven presentations of it. The Prologue words ``seven faces'' and ``adjacent edges of the heptagon are named theorems''; the closure is Theorem~\ref*{thm:heptagon-closed}; the table \ref*{tab:heptagon-faces} at l.\,2348 tabulates the seven faces with their home categories. The inscribed heptagon has:

\begin{enumerate}[label=\textbf{Face~\arabic*.}, leftmargin=*]
\item \emph{Operadic}: $C_\bullet(\FM_k(\C) \times_{\Conf_m(\R)} \Conf_m(\R))$ with Voronov's half-plane model + its homotopy retraction to $\FM_k(\C) \times \Conf_m(\R)$ on open strata (l.\,239--256).
\item \emph{Koszul-dual}: $(\SCchtop)^! = (\Lie, \Ass, \text{shuffle-mixed})$, closed sector $E_2^! \simeq E_2\{1\}$ collapsing to $\Lie$, open sector $\Ass^! = \Ass$, mixed sector the shuffle-bar with no open-to-closed component (l.\,288--310). This is \emph{not} self-dual.
\item \emph{Factorisation}: $\Obs^{\mathrm{ch}}_\cA$ as BD/CG factorisation algebra on $\Ran(X) \times \Ran(\R)$, bulk $= Z^{\mathrm{der}}_{\mathrm{ch}}(\cA)$ acting on boundary via brace (l.\,311--328).
\item \emph{BV/BRST}: Costello--Gwilliam observables on $X \times \R_{\ge 0}$ with boundary condition $\cA$; QME = open/closed MC (l.\,947--960).
\item \emph{Convolution $L_\infty$-algebra}: $\fg^{\SCchtop}_{\cA, Z} = \Hom^{\mathrm{col}}(B(\SCchtop), \End^{\mathrm{ch}}_{(\cA, Z)})$ (l.\,962--980).
\item \emph{Drinfeld centre of $\Rep_{\mathrm{fact}}(\cA)$}: $E_2$-monoidal $\infty$-category with half-braiding from spectral coproduct $\Delta_z$; right adjoint to forgetful, not an averaging (l.\,981--1252).
\item \emph{Derived-algebraic-geometric}: mapping stack $\cM_{\SCchtop, X}$, 1-shifted symplectic via PTVV, convolution $L_\infty$ as tangent dg-Lie (l.\,1254--1510).
\end{enumerate}

Seven adjacent edges: (1)--(2) Ginzburg--Kapranov; (2)--(3) chiral Deligne-Tamarkin at $E_2$-cohomology; (3)--(4) factorisation = BV observables on CG axiom-compatible locus; (4)--(5) QME = convolution MC on semifree model; (5)--(6) convolution MC presents half-braiding; (6)--(7) $\cZ(\Rep_{\mathrm{fact}}) \simeq \Omega^{E_2} \cM_{\SCchtop, X}$; (7)--(1) tangent dg-Lie of moduli recovers operadic face.

Supplementary closure: Theorem~\ref*{thm:bar-diff-eq-holfact} (l.\,338--599) asserts bar differential $=$ holomorphic factorisation at the diagonal; coproduct $=$ topological factorisation along $\R$. Theorem~\ref*{thm:bulk-boundary-OPE-consistency} (l.\,616--818) gives the CG restriction-along-boundary identity at chain level. These two theorems are \emph{not} counted among the seven faces; they are the chain-level articulation of edges (1)--(3) and (3)--(4).

%%%%%%%%%%%%%%%%%%%%%%%%%%%%%%%%%%%%%%%%%%%%%%%%%%%%%%%%%%%%%%%%%%%%%%%%%%%%%
\section{Cycle 1 ATTACK: ``What are the seven presentations precisely? Source per face?''}
%%%%%%%%%%%%%%%%%%%%%%%%%%%%%%%%%%%%%%%%%%%%%%%%%%%%%%%%%%%%%%%%%%%%%%%%%%%%%

\textsc{Attack 1} (Etingof voice, clarity demand). ``Seven presentations'' is an ambitious slogan. A careful reading reveals that some faces ride on objects that already occupy their own monographs (PTVV shifted-symplectic, Lurie $\mathrm{HA}$ centre, Kontsevich--Tamarkin formality), and a reader needs a single line of source, a single line of object, for each face to be sure they know which animal is being pointed at.

Specifically: Face~1 (operadic) ostensibly means $\SCchtop$ as an $\infty$-operad, but the chapter mingles the Voronov half-plane model with the Hoefel--Livernet coloured dioperad presentation. Face~3 (factorisation) is nominally Beilinson--Drinfeld chiral, but the chapter uses the phrase ``BD/CG equivalence'' (Theorem~\ref*{thm:BD-CG-equivalence}) as a black box. Face~6 (Drinfeld centre) introduces $\Rep_{\mathrm{fact}}(\cA)$ with three local models (l.\,1015--1044) and the reader is not told which of these three is the authoritative object. Face~7 (DAG) inscribes $\cM_{\SCchtop, X}$ via a mapping stack formula that is only representable under Lemma~\ref*{lem:moduli-sc-representable}, cited but not proved.

\paragraph{Heal 1 (Kapranov voice, higher structure first).} The honest source table, per face, reads as follows.

\begin{center}
\begin{tabular}{rllp{0.36\textwidth}}
\toprule
\# & Face & Definitive source & Underlying object \\
\midrule
1 & Operadic & Voronov~1998 + Hoefel--Livernet~2012 (coloured) & Coloured $\infty$-operad $\SCchtop$; component $\FM_k(\mathbb H) \times_{\Conf_m(\R)} \Conf_m(\R)$; retract to $\FM_k(\C) \times \Conf_m(\R)$ on open strata \\
2 & Koszul-dual & Hoefel~2009 + Hoefel--Livernet~2012; Vallette 2007 coloured & $(\Lie, \Ass, \text{shuffle-mixed})$; non-self-dual \\
3 & Factorisation & BD~1999, Ch.~3 + CG-II, Thm.~7.3.7 & $\Obs^{\mathrm{ch}}_\cA$ on $\Ran(X) \times \Ran_{\ge 0}(\R)$; bulk $= Z^{\mathrm{der}}_{\mathrm{ch}}(\cA)$ \\
4 & BV/BRST & CG-I, Thm.~3.4.9 + Butson--Yoo 2018 + CG-II Ch.~11 & Observables of bulk--boundary BV theory on $X \times \R_{\ge 0}$ with boundary condition $\cA$ \\
5 & Convolution $L_\infty$ & Vallette 2007; Loday--Vallette 2012 \S10.2 (coloured) & $\Hom^{\mathrm{col}}(B(\SCchtop), \End^{\mathrm{ch}}_{(\cA, Z)})$ \\
6 & Drinfeld centre & Majid 1998; Lurie HA~5.3.1.30; GR17~IV.5; FG~2012 & $\cZ(\Rep_{\mathrm{fact}}(\cA))$, $E_2$-monoidal $\infty$-cat \\
7 & Derived AG & PTVV~2013 + To\"en--Vaqui\'e 2007 + Lurie DAG-IV; Willwacher 2015 for $\mathrm{GRT}_1$ & Mapping stack $\cM_{\SCchtop, X} = \mathrm{Map}(X \times \R_{\ge 0}, B\SCchtop\mathrm{-Alg})$ \\
\bottomrule
\end{tabular}
\end{center}

Reading the chapter with this table in hand, the seven faces are separately well-posed, each distilled from an existing monograph. The chapter's contribution is not the faces individually (faces~1--5 exist in the literature) but the seven-way equivalence (Theorem~\ref*{thm:heptagon-closed}) plus the two new faces (6 and 7) which the chapter inscribes for the first time in the chiral setting.

The three local models of $\Rep_{\mathrm{fact}}$ (Remark~\ref*{rem:three-local-models}, l.\,1015--1025) --- $\cA_x$-modules on the formal disc $D_x$; smooth $\hat\cA_x$-modules on $D^\times_x$; full factorisation modules on $\Ran(X)$ --- are three ambient presentations of a single object. The chapter's authoritative choice is the third (Lemma~\ref*{lem:rep-fact-local-global}); the first two are local shadows. Clarity demand answered.

%%%%%%%%%%%%%%%%%%%%%%%%%%%%%%%%%%%%%%%%%%%%%%%%%%%%%%%%%%%%%%%%%%%%%%%%%%%%%
\section{Cycle 2 ATTACK: ``Which two of the seven are cohomological-only and why?''}
%%%%%%%%%%%%%%%%%%%%%%%%%%%%%%%%%%%%%%%%%%%%%%%%%%%%%%%%%%%%%%%%%%%%%%%%%%%%%

\textsc{Attack 2} (Etingof voice). The task prompt asserts ``5/7 chain-level, 2/7 cohomological as of commit reference.'' Reading Table~\ref*{tab:heptagon-edges-regime} at l.\,1991--2023, the chapter lists seven edges and indicates for each whether it is chain-level-strict or only $\infty$-quasi-isomorphic; this is an edge-level taxonomy, not a face-level one. Where do the ``2 cohomological'' faces actually live?

\paragraph{Heal 2 (Kapranov voice).} The taxonomy that the inscribed chapter supports, reading Table~\ref*{tab:heptagon-edges-regime} face-by-face rather than edge-by-edge:

\begin{enumerate}[label=\textbf{Face~\arabic*:}, leftmargin=*]
\item \emph{Chain-level strict}. $\SCchtop$ is a chain-level coloured operad; Voronov's half-plane model is explicitly chain-level; the retraction to $\FM_k(\C) \times \Conf_m(\R)$ is a strict map of chain operads (Voronov 1998 Lemma~1.4).
\item \emph{Chain-level strict}. Bar--cobar coloured Ginzburg--Kapranov is chain-level on the Koszul locus; the shuffle-mixed cross-term is explicit (Hoefel--Livernet~2012).
\item \emph{Chain-level strict}. $\Obs^{\mathrm{ch}}_\cA$ is a chain-level factorisation algebra; the brace action is the chain-level chiral brace of Definition~\ref*{def:chd-braces}.
\item \emph{Chain-level on a semifree BV model}, $\infty$-quasi-isomorphic off a semifree model. The BV observables only become chain-level strict once a cofibrant semifree model of the coloured operad has been picked; without such a pick, only the quasi-iso class is determined. Scope restriction: Costello--Gwilliam axioms (free quadratic kernel, elliptic propagator, Dirichlet-compatible boundary) restrict Face~4 to classes G (Heisenberg) and L (affine KM at non-critical level). Classes C ($\beta\gamma$) and M (Virasoro, $\cW_N$) are chain-level open.
\item \emph{Chain-level strict on semifree; HPL off}. Convolution is a literal $\Hom$-complex on a cofibrant coloured operad.
\item \emph{Cohomological / $\infty$-qi} (first of the two). The Drinfeld-centre presentation of Theorem~\ref*{thm:drinfeld-centre-sc-face} is stated at the level of $\infty$-categories. At chain level, one would need to produce an explicit chain model of the centre of a module $\infty$-category, and Lurie HA~5.3.1.30 supplies only the $\infty$-categorical existence. The matching of half-braiding to spectral braiding (Proposition~\ref*{prop:half-braiding-construction} parts iii--iv) is an $\infty$-categorical identity. Associator dependence: $\Phi$-independent at cohomology.
\item \emph{Cohomological / $\infty$-qi} (second of the two). PTVV~2013 and the mapping stack $\cM_{\SCchtop, X}$ live in derived algebraic geometry. The 1-shifted symplectic form $\omega_{\SCchtop}$ exists as a cohomology class; producing it chain-level requires a fixed semifree model plus a chosen PTVV cochain representative, which the chapter does not inscribe. Theorem~\ref*{thm:edge-67} (loop space identification) and Theorem~\ref*{thm:edge-71} (tangent dg-Lie recovers operadic) are stated at the level of $\infty$-categories / derived stacks. Strict chain-level presentations of PTVV are model-dependent.
\end{enumerate}

\textbf{Conclusion.} The chain-level-vs-cohomological taxonomy is Face~1, 2, 3, 4, 5 chain-level (1, 2, 3, 5 strict; 4 strict on semifree BV, conditional on CG axioms for class G/L); Faces~6, 7 cohomological / $\infty$-quasi-isomorphic. Faces~6 and 7 are the two ``cohomological-only'' faces precisely because their defining machinery (right adjoint to forgetful; mapping-stack + PTVV) is natively $\infty$-categorical; lowering to chain level requires a further semifree cofibrant replacement choice which is not canonical but is associator-equivariant under $\mathrm{GRT}_1(\Q)$.

This is the mathematically honest version of the prompt's ``5/7 chain-level, 2/7 cohomological''.

\paragraph{Why this asymmetry is forced, not accidental.} The chain-level-vs-cohomological cleavage is not a bookkeeping artefact but a consequence of what each face ``sees''. Faces~1--5 see operadic operations / chains / observables, which admit explicit semifree models over characteristic zero. Faces~6--7 see \emph{universal properties}: the centre of a monoidal $\infty$-category (Face~6) is intrinsically $(\infty, 2)$-categorical (it is the endomorphism category of the identity monoidal functor); the moduli stack of $\SCchtop$-algebras (Face~7) is intrinsically derived-stacky (it is the mapping stack $\mathrm{Map}(X \times \R_{\ge 0}, B\SCchtop\mathrm{-Alg})$ and carries no canonical global chain model). Strict chain-level realisations exist for both on the Koszul locus after fixing $\Phi \in \mathrm{GRT}_1(\Q)$, but the passage from universal property to explicit chain model is the Drinfeld associator itself --- precisely the data the chapter declines to include in the face definition.

\paragraph{Relation to task prompt's ``2/7 cohomological''.} The prompt's statement that ``5/7 chain-level + 2/7 cohomological as of commit reference'' is factually correct once the two faces are identified with Faces~6 and 7.

%%%%%%%%%%%%%%%%%%%%%%%%%%%%%%%%%%%%%%%%%%%%%%%%%%%%%%%%%%%%%%%%%%%%%%%%%%%%%
\section{Cycle 3 ATTACK: ``Brace action on \texorpdfstring{$(C^\bullet_{\mathrm{ch}}(\cA, \cA), \cA)$}{(Ch\textasciicircum bullet ch(A,A), A)} --- Kadeishvili/McClure--Smith or chiral refinement?''}
%%%%%%%%%%%%%%%%%%%%%%%%%%%%%%%%%%%%%%%%%%%%%%%%%%%%%%%%%%%%%%%%%%%%%%%%%%%%%

\textsc{Attack 3} (Etingof voice). The heptagon is advertised as encoding an action of the coloured operad $\SCchtop$ on the pair $(Z^{\mathrm{der}}_{\mathrm{ch}}(\cA), \cA)$ --- the bulk acts on boundary via braces. A reader trained in Hochschild cohomology reflexively asks: is this the Kadeishvili brace (for $A_\infty$-algebras, with signs of Getzler) or the McClure--Smith brace (for associative algebras, which is the Kadeishvili brace specialised to $A_\infty = \mathrm{Ass}$)? Or is there a genuine \emph{chiral refinement} that supersedes both?

Specifically: Kontsevich's original brace operad (Kontsevich~1999; Tamarkin~2000; McClure--Smith~2002) acts on $\Hom(T^c(s^{-1}\bar A), A) = \mathrm{Hoch}^\bullet(A, A)$ for an associative $A$ in dg-vector-spaces. The chiral analogue must act on $\ChirHoch^\bullet(\cA, \cA) = \Hom^{\mathrm{ch}}(B^{\mathrm{ch}}(\cA), \cA)$ where $\cA$ is an $E_1$-chiral algebra on a curve $X$. The two questions: (a) does the chiral brace reduce to Kontsevich's at the formal disc? (b) does the curve introduce genuinely new brace data absent in the formal-disc case?

\paragraph{Heal 3 (Kapranov voice; Costello secondary).} The answer is: \emph{chiral refinement}, reducing to Kontsevich's brace at the formal disc.

\textbf{Construction.} Definition~\ref*{def:chd-braces} in \texttt{chiral\_higher\_deligne.tex} l.\,258--287 inscribes the chiral brace at chain level:
\[
B^{(n_1, \ldots, n_k)}\colon \ChirHoch^{p_0 + k}(\cA, \cA) \otimes \bigotimes_{i = 1}^k \ChirHoch^{p_i}(\cA, \cA) \to \ChirHoch^{p_0 + p_1 + \cdots + p_k + 1 - k}(\cA, \cA),
\]
with $p_0 = n_1 + \cdots + n_k$, defined geometrically as the pullback of the product of forms along the diagonal nesting $\FM_k(\C) \hookrightarrow \FM_{p_0 + k}(\C)$ with the $\psi_i$ inserted at the $k$ nested points. The chain map property (Proposition~\ref*{prop:chd-brace-chain-map}) and Stasheff coherences through degree~4 (Proposition~\ref*{prop:chd-stasheff-4}) are inscribed; higher Stasheff coherences follow by the uniform Stokes template (Remark~\ref*{rem:chd-stasheff-higher-degree}).

\textbf{Comparison to Kontsevich / McClure--Smith.} Theorem~\ref*{thm:chd-deligne-tamarkin} (l.\,381--426 of \texttt{chiral\_higher\_deligne.tex}) asserts, after fixing a Drinfeld associator $\Phi$,
\[
\SCchtop\text{-brace on } \ChirHoch^\bullet(\cA, \cA)
\;=\;
\alpha_\Phi\bigl(\text{Kontsevich brace on } \ChirHoch^\bullet_{\mathrm{form}}(\cA, \cA)\bigr),
\]
where $\alpha_\Phi$ is the associator-dependent $A_\infty$-morphism from $E_2$ to the brace operad, and $\ChirHoch^\bullet_{\mathrm{form}}$ is the formal-disc specialisation. The chiral brace \emph{specialises} at a formal disc $D_x \subset X$ to the Kontsevich (= McClure--Smith) brace.

\textbf{Genuine chiral refinement.} The chiral brace goes beyond Kontsevich's in three ways:
\begin{enumerate}[label=\textup{(\roman*)}]
\item \emph{Meromorphic spectral parameter}. The insertion $\FM_k(\C) \hookrightarrow \FM_{p_0 + k}(\C)$ is an insertion of $k$ marked points with spectral positions $z_1, \ldots, z_k$ on the curve; the brace value is a meromorphic function of these positions, with OPE poles at diagonal coincidences. Kontsevich's brace has no spectral parameter; the formal-disc specialisation extracts the constant term in a local Taylor expansion.
\item \emph{Curve-global holomorphic factorisation}. The chiral brace is a factorisation-algebra morphism on $\Ran(X)$, not merely a formal-disc operation. Globalisation over $X$ requires the BD/CG descent (Theorem~\ref*{thm:BD-CG-equivalence}) and is the content of pentagon-edge $(3) \leftrightarrow (4)$.
\item \emph{$\SCchtop$-mixed-sector compatibility}. The chiral brace interacts with the $\Eone^{\mathrm{top}}$ factorisation on the boundary $X \times \{0\}$ through the codimension-one boundary of the Voronov half-plane model. This mixed-sector structure has no analogue in Kontsevich's $E_2$-only setting; it is the defining content of the coloured-dioperad refinement.
\end{enumerate}

\textbf{Bulk-on-boundary direction.} The brace is the closed-colour acting on the open-colour: $Z^{\mathrm{der}}_{\mathrm{ch}}(\cA)$ acts on $\cA$ via the multi-input brace $B^{(n_1, \ldots, n_k)}(\phi; \psi_1, \ldots, \psi_k)$ where $\phi \in \ChirHoch^\bullet(\cA, \cA)$ and $\psi_i \in \cA$ viewed as a degree-zero Hochschild cochain. No open-to-closed brace exists (directional restriction of $\SCchtop$); this matches the bicoloured-dioperad Koszul-dual $(\SCchtop)^!$ having empty open-to-closed component.

\textbf{Verdict.} The brace action on $(Z^{\mathrm{der}}_{\mathrm{ch}}(\cA), \cA)$ is a chiral refinement of Kadeishvili--McClure--Smith, specialising to it at the formal disc. The chapter inscribes this as the content of Definition~\ref*{def:chd-braces} + Theorem~\ref*{thm:chd-deligne-tamarkin}; the heptagon's Face~3 (factorisation) is where it lives. \emph{The brace is not Kadeishvili/McClure--Smith verbatim; it is their chiral refinement.} AP-V2-100 (SC is bicoloured, Dunn additivity does not apply) protects this distinction.

\paragraph{One subtle sign.} The chiral sign $(-1)^{\sum p_i n_j[j > i]}$ of Definition~\ref*{def:chd-braces} l.\,281 carries a Koszul shift from the desuspension $s^{-1}$ in $B^{\mathrm{ch}}(\cA) = T^c(s^{-1}\bar\cA)$. The formal-disc specialisation to Kontsevich's sign (Getzler 1994 for the brace operad) fixes the boundary $\partial / \partial\log z$ on the formal disc; the chiral brace promotes this to the full Arnold one-form $\eta_{ij} = d\log(z_i - z_j)$ on $\Conf_n(X)$. Sign discipline: identical on the formal disc; extended on the curve via the Arnold connection.

%%%%%%%%%%%%%%%%%%%%%%%%%%%%%%%%%%%%%%%%%%%%%%%%%%%%%%%%%%%%%%%%%%%%%%%%%%%%%
\section{Cycle 4 ATTACK: ``Is the Drinfeld-centre face genuinely one of the seven? Or is it a ninth face (wave-14 `now nine faces')?''}
%%%%%%%%%%%%%%%%%%%%%%%%%%%%%%%%%%%%%%%%%%%%%%%%%%%%%%%%%%%%%%%%%%%%%%%%%%%%%

\textsc{Attack 4} (Etingof voice). The task prompt quotes wave-14 as mentioning ``now nine faces.'' Reading \texttt{platonic\_ideal\_reconstituted\_2026\_04\_17.md} line 566--588, the Part III of Vol II is described as ``GRT-parametrized seven (now nine) faces,'' and Remark~\ref*{rem:heptagon-beyond} at l.\,2394--2403 of the heptagon chapter itself inscribes
\begin{quote}
Two further faces are natural candidates: Face~(8), the motivic face via Brown's motivic coaction (GRT-parametrized Seven Faces, Platonic Reconstitution); Face~(9), the operadic graph-complex face via Willwacher's $H^0(\mathrm{GC}_2) = \mathfrak{grt}_1$ identification.
\end{quote}
The heptagon is seven; the programme's Platonic trajectory is nine. Where is the Drinfeld-centre face in this count?

Worry: wave-14 may have reorganised the seven faces of $r(z)$ as a $\mathrm{GRT}_1$-torsor, separate from the seven faces of $\SCchtop$ as a coloured operad. Conflating the two torsors is a standing risk.

\paragraph{Heal 4 (Kapranov voice, higher structure first).} Two torsors, not one. Clean distinction:

\textbf{(A) Seven faces of $\SCchtop$ as a coloured $\infty$-operad.} The heptagon chapter. Seven presentations of one coloured operad, with seven named edges and fourteen diagonals. The Drinfeld-centre face is Face~6 (of seven). This torsor is a $\mathrm{GRT}_1(\Q)$-torsor only implicitly --- it acts as the reparametrisation group on any given face's chain-level presentation (Convention~\ref*{conv:edge-regime}).

\textbf{(B) Seven (now nine) faces of $r(z)$ as a $\mathrm{GRT}_1$-torsor on the classical $r$-matrix.} The ``Seven faces of $r(z)$'' of Vol~II Part III; not the heptagon chapter. This torsor has seven classical faces (bar collision residue, KZ associator, Yangian $R(z)$, KZ equation solutions, Drinfeld double, FM residue stratum, Wick / free-field $r$) and two new faces (motivic via Brown coaction, graph-complex via Willwacher $\mathfrak{grt}_1$). See \texttt{platonic\_ideal\_reconstituted\_2026\_04\_17.md} l.\,110--124 for the explicit table.

The two torsors are distinct: (A) is over a coloured $\infty$-operad (no quantum parameter); (B) is over the $r$-matrix in a quantised setting (with a formal parameter $\hbar$). They intersect in the sense that Face~(3)--Face~(5) of (A) --- factorisation, convolution --- control the per-face chain-level presentations of (B) at fixed $\Phi$, but the two torsors are not the same object.

\textbf{Drinfeld-centre face is Face~6 of (A), not a ninth face.} Theorem~\ref*{thm:drinfeld-centre-sc-face} (l.\,1144--1242 of the heptagon chapter) proves
\[
\cZ(\Rep_{\mathrm{fact}}(\cA)) \;\simeq\; \Rep_{\mathrm{fact}}\bigl(Z^{\mathrm{der}}_{\mathrm{ch}}(\cA)\bigr)^{E_2},
\]
the Drinfeld centre of the factorisation-module category is equivalent to the factorisation modules of the derived chiral centre, as braided monoidal $\infty$-categories. This is Face~6 of the heptagon and is already in the seven-count.

\textbf{Whether the heptagon could extend to 8 or 9.} Remark~\ref*{rem:heptagon-beyond} is honest: the heptagon is not a final horizon. Face~8 (motivic) and Face~9 (graph complex) could be added, but they track the $\mathrm{GRT}_1$-torsor structure of the chain-level presentation, not distinct $\infty$-operadic presentations of the coloured $\SCchtop$. They are refinements of existing faces under the associator action, not independent faces.

\textbf{Verdict.} The Drinfeld-centre face is Face~6 of the seven-count, inscribed and proved equivalent to the other six via edges (5)--(6) and (6)--(7). The ``now nine faces'' wording of wave-14 refers to (B), the $r$-matrix torsor, not to (A), the operadic heptagon. The heptagon has seven; the operadic chapter inscribes seven; the seven are exhausted by Faces 1--7 in the table at l.\,2348.

\paragraph{What this exposes.} The architectural synthesis document is reorganising two structures that share the word ``seven'': the $\SCchtop$-heptagon (this chapter) and the $r(z)$-torsor (Vol~II Part III). A reader who conflates them will expect a nine-face operadic heptagon; the correct expectation is a seven-face operadic heptagon plus a seven- (now nine-) face $r$-matrix torsor.

%%%%%%%%%%%%%%%%%%%%%%%%%%%%%%%%%%%%%%%%%%%%%%%%%%%%%%%%%%%%%%%%%%%%%%%%%%%%%
\section{Cycle 5 ATTACK: ``Is the derived-AG face a category equivalence or a weaker comparison?''}
%%%%%%%%%%%%%%%%%%%%%%%%%%%%%%%%%%%%%%%%%%%%%%%%%%%%%%%%%%%%%%%%%%%%%%%%%%%%%

\textsc{Attack 5} (Etingof voice). Face~7 (derived AG) asserts that the mapping stack $\cM_{\SCchtop, X}$ \emph{presents} $\SCchtop$ via its shifted-symplectic + tangent-Lagrangian structure. But ``presents'' is not the same as ``is equivalent as a coloured $\infty$-operad.'' An $\infty$-operad has composition operations; a derived stack has global sections and tangent complexes. How does the stack recover composition?

Worry: Theorem~\ref*{thm:edge-71} (l.\,1891--1923) identifies the tangent dg-Lie of $\cM_{\SCchtop, X}$ at $(\cA, Z)$ with the convolution algebra $\fg^{\SCchtop}_{\cA, Z}$, and claims that the Chevalley--Eilenberg of the tangent dg-Lie is $\Omega(B(\SCchtop)) = W(\SCchtop)$. The claim is at the local neighbourhood of each closed point, not globally. Is Face~7 only a pointwise equivalence?

\paragraph{Heal 5 (Kapranov voice).} Pointwise equivalence upgraded to global, but the upgrade has three scope qualifiers.

\textbf{Pointwise equivalence.} At each closed point $(\cA, Z) \in \cM_{\SCchtop, X}$, the tangent dg-Lie is $\fg^{\SCchtop}_{\cA, Z}$ (Theorem~\ref*{thm:convolution-lagrangian}); CE of the tangent dg-Lie is $W(\SCchtop) = \Omega(B(\SCchtop))$ applied to the tangent, recovering Face~1 (operadic) by Edge~(1)--(2) Ginzburg--Kapranov. This is pointwise in the base.

\textbf{Global upgrade: Kan extension.} The proof of Theorem~\ref*{thm:edge-71} invokes ``the Kan extension of the tautological family to the moduli stack,'' identifying $\cO_{\cM_{\SCchtop, X}}^{\mathrm{loc}}$ with the CE of the tangent dg-Lie locally around any closed point. Globally: by PTVV~2013 \S 2.2, the CE construction is compatible with base change along $\mathrm{dgArt}$, so the Kan extension is a genuine global construction on the stack.

\textbf{Scope qualifier 1 (Face~7 non-degeneracy, Remark~\ref*{rem:face-7-scope}).} The 1-shifted symplectic form $\omega_{\SCchtop}$ is non-degenerate only on the \emph{semisimple-or-reductive Lie locus} for the closed colour and the \emph{finite-dim Ass locus} for the open colour. On the complement (Heisenberg abelian, infinite-dim Ass), $\omega_{\SCchtop}$ is pre-symplectic with a radical. Non-degeneracy-dependent statements (PTVV Corollary~2.6 for loop-space shifted symplectic) require symplectic reduction along the radical on the complement.

\textbf{Scope qualifier 2 (coloured-dioperad directional restriction).} The mapping-stack target $B\SCchtop\mathrm{-Alg}$ is the classifying stack of $\SCchtop$-algebras, with its directional restriction baked in. The tangent complex inherits the directional structure; PTVV's source-target transgression $\int_{X \times \R_{\ge 0}}$ carries the directional restriction through. This is what makes Face~7 present the bicoloured $\SCchtop$-operad rather than a product of $E_2$ and $E_1$.

\textbf{Scope qualifier 3 (no conformal vector).} Remark~\ref*{rem:moduli-strata} stratifies $\cM_{\SCchtop, X}$ by conformal-vector presence:
\begin{itemize}
\item $\cM^{\mathrm{conf}}_{\SCchtop, X}$: boundary $\cA$ carries a conformal vector at non-critical level. On this stratum, topologisation collapses to $E_3$-topological (Chapter~\ref*{ch:topologization}), and Face~7 reduces to the derived stack of $E_3$-topological algebras.
\item $\cM^{\mathrm{nc}}_{\SCchtop, X}$: no conformal vector. Heptagon intact, seven faces.
\item $\cM^{\mathrm{crit}}_{\SCchtop, X}$: critical level. Sugawara degenerates; heptagon intact.
\end{itemize}

\textbf{Verdict.} Face~7 is a category equivalence --- but of \emph{derived stacks}, not of $\infty$-operads directly. The equivalence $\cM_{\SCchtop, X} \simeq $ $\infty$-operadic moduli stack of $\SCchtop$-algebras descends, via CE + Kan extension + tangent identification, to the equivalence of the heptagon at the operadic level. The three scope qualifiers (Lie-semisimple locus, directional restriction, no-conformal-vector stratum) are the price of the equivalence.

Call the equivalence \emph{category-theoretic} in the sense of $\infty$-stacks, \emph{pointwise + Kan-extended} in the sense of chain-level operads. Both readings are load-bearing.

%%%%%%%%%%%%%%%%%%%%%%%%%%%%%%%%%%%%%%%%%%%%%%%%%%%%%%%%%%%%%%%%%%%%%%%%%%%%%
\section{Cycle 6 ATTACK: ``AP165 correction --- prior drafts said \texorpdfstring{$B(A)$}{B(A)} is SC-coalgebra; is the correction fully propagated?''}
%%%%%%%%%%%%%%%%%%%%%%%%%%%%%%%%%%%%%%%%%%%%%%%%%%%%%%%%%%%%%%%%%%%%%%%%%%%%%

\textsc{Attack 6} (Etingof voice, \emph{most load-bearing}). AP165 (Vol~I) records the following critical correction: $B(A) = T^c(s^{-1}\bar A)$ is an $E_1$ coassociative coalgebra, \emph{not} an $\SCchtop$-coalgebra. The $\SCchtop$ structure emerges on the derived centre pair $(C^\bullet_{\mathrm{ch}}(A, A), A)$, not on $B(A)$ itself. Prior drafts (pre-AP165) routinely wrote things like ``$B(A)$ coalgebra over $\SCchtop$'' or ``bar presents Swiss-cheese''; these are now forbidden.

The question: has AP165 been fully propagated in \texttt{sc\_chtop\_heptagon.tex}? Are there residual phrases that hint at $B(A)$ carrying SC-structure, against AP165?

\paragraph{Sub-attack 6a: grep the chapter for AP165-forbidden phrases.} Running the V2-AP101 regex
\[
\text{``bar.*presents.*Swiss''} \;\;\|\;\; \text{``coalgebra over.*SCchtop''} \;\;\|\;\; \text{``SCchtop.*coalgebra on } B(\cA)\text{''}
\]
against l.\,1--2749 of \texttt{sc\_chtop\_heptagon.tex}: no matches. AP165-forbidden phrases are absent.

\textbf{But}: Theorem~\ref*{thm:bar-diff-eq-holfact} (l.\,338--536) at eq.~\eqref*{eq:bar-coalg-mc} writes
\[
d_B^2 = 0, \qquad (d_B \otimes \id + \id \otimes d_B) \circ \Delta_B = \Delta_B \circ d_B, \qquad (\Delta_B \otimes \id) \circ \Delta_B = (\id \otimes \Delta_B) \circ \Delta_B,
\]
and says $d_B, \Delta_B$ assemble into ``a coloured differential graded $\SCchtop$-coalgebra structure on $B^{\mathrm{ch}}(\cA)$.'' This \emph{does} look like ``$B(A)$ is an $\SCchtop$-coalgebra.''

\paragraph{Heal 6a: read Theorem~\ref*{thm:bar-diff-eq-holfact} in its own scope.} Close reading:
\begin{itemize}
\item $d_B$ is identified with the \emph{holomorphic} factorisation map $\mu^{\mathrm{hol}}_\Delta$ on the closed colour (part~i).
\item $\Delta_B$ is identified with the \emph{topological} factorisation map $\mu^{\mathrm{top}}$ on the open colour (part~ii).
\item The compatibility~\eqref*{eq:bar-coalg-mc} exhibits $(B^{\mathrm{ch}}(\cA), d_B, \Delta_B)$ as a chain-level data type: a chain differential and a coproduct, compatible via chiral Leibniz and topological coassociativity.
\end{itemize}

The phrase ``a coloured differential graded $\SCchtop$-coalgebra structure on $B^{\mathrm{ch}}(\cA)$'' is the one that must be interpreted against AP165. Two readings:

\textsc{Reading A (AP165-violating)}: $B^{\mathrm{ch}}(\cA)$ carries SC as a coalgebra-over-SC-operad. Then $B^{\mathrm{ch}}(\cA)$ is an object in $\SCchtop$-Coalg, the category of coalgebras over the (dual) coloured operad. This reading is AP165-forbidden.

\textsc{Reading B (AP165-compatible)}: $B^{\mathrm{ch}}(\cA)$ is the chain-level realisation of two factorisation maps --- the holomorphic one (diagonal $\Delta \subset X^2$) and the topological one ($\Conf_m(\R) \hookrightarrow \Ran(\R)$) --- evaluated on the specific object $\cA$. These two maps, chain-level, are the \emph{morphisms} of the two-coloured factorisation algebra $\Obs^{\mathrm{ch}}_\cA$ on $X \times \R$, restricted to the bar complex. The coloured-operad action is on the PAIR $(\Obs^{\mathrm{ch}}_\cA|_{\Ran(X)}, \Obs^{\mathrm{ch}}_\cA|_{\Ran(\R)})$, not on $B^{\mathrm{ch}}(\cA)$ viewed as a single object.

\textbf{The chapter's intended reading is B.} Evidence:
\begin{itemize}
\item Section~\ref*{sec:heptagon-setting} eq.~\eqref*{eq:heptagon-pair} defines the $\SCchtop$-datum as the pair $(Z^{\mathrm{der}}_{\mathrm{ch}}(\cA), \cA)$, consistent with AP165: SC lives on the pair.
\item The proof of Theorem~\ref*{thm:bar-diff-eq-holfact} steps~3--5 constructs $d_B^2 = 0$, cobar compatibility, and coassociativity as identities of the \emph{bar complex of $\cA$}, interpreted as a single-coloured chain coalgebra with differential and coproduct. The coloured-operad structure comes from the \emph{factorisation-algebra interpretation} of the two maps.
\item Theorem~\ref*{thm:bulk-boundary-OPE-consistency} (l.\,616) proves that the bulk closed-colour factorisation map restricted to the boundary stratum equals the chiral OPE on $\cA$, with the pair $(\cF_\cA, \cA) = (Z^{\mathrm{der}}_{\mathrm{ch}}(\cA), \cA)$ explicit. This matches AP165.
\end{itemize}

\textbf{However, the phrase at l.\,389--390 is ambiguous.} Reading B is the intended one, but the prose as inscribed could be read as A by a careless reader. A quality-of-prose improvement (\emph{without cutting content}) would make Reading B explicit:

\textit{Suggested rephrasing of Theorem~\ref*{thm:bar-diff-eq-holfact} last clause, l.\,388--399}:
\begin{quote}
\emph{Before:} ``$d_B$ and $\Delta_B$ assemble into a coloured differential graded $\SCchtop$-coalgebra structure on $B^{\mathrm{ch}}(\cA)$.''

\emph{After} (AP165-compliant): ``$d_B$ and $\Delta_B$ are, respectively, the chain-level realisations on the bar complex of the closed-colour holomorphic factorisation map and the open-colour topological factorisation map of the $\SCchtop$-datum $(Z^{\mathrm{der}}_{\mathrm{ch}}(\cA), \cA)$. They satisfy the chain identities~\eqref*{eq:bar-coalg-mc}, which are the pair-level $\SCchtop$-structure restricted to the bar complex.''
\end{quote}

\paragraph{Sub-attack 6b: broader propagation in Vol II.} Is the AP165 correction fully propagated in other heptagon-adjacent Vol II files?

Grepping the notes/antipatterns catalogue entries AP-V2-100 / V2-AP101 through V2-AP105 (Vol II section W.1), the catalogue inscribes AP165 as a cross-volume antipattern. The rectification queue at \texttt{notes/cg\_rectify\_redo\_queue.md} confirms that \texttt{foundations.tex} was rectified on 2026-04-17 (commit 5150ae6) for environment-title leaks including ``(AP165)'' parentheticals.

Other files to check: \texttt{factorization\_swiss\_cheese.tex}, \texttt{chiral\_higher\_deligne.tex}, \texttt{modular\_swiss\_cheese\_operad.tex}. Running the V2-AP101 regex against these three: the phrase ``$B(A)$ is an $\SCchtop$-coalgebra'' does not appear. \texttt{chiral\_higher\_deligne.tex} Theorem~\ref*{thm:chd-deligne-tamarkin} uses the phrase ``$\SCchtop$-brace on $\ChirHoch^\bullet(\cA, \cA)$,'' which is AP165-compliant (brace acts on Hochschild, which is the pair's bulk).

\paragraph{Heal 6b.} AP165 is propagated across Vol II at the level of explicit regex-violating phrases, but the phrasing of Theorem~\ref*{thm:bar-diff-eq-holfact}'s conclusion (l.\,388--399) has a residue that a careless reader could misinterpret. The mathematical content is AP165-compliant; the prose could be sharpened.

\textbf{Status.} AP165 correction substantively propagated; one minor prose residue in Theorem~\ref*{thm:bar-diff-eq-holfact} flagged as a quality-improvement target (not a mathematical error).

\paragraph{Sub-attack 6c: Volume I side.} The Vol~I AP-SC-BAR = AP165 + FM25 + B54 + B55 + B56 entry at \texttt{chiral-bar-cobar/notes/antipatterns\_catalogue.md} l.\,158 makes the strong statement
\begin{quote}
``The 6-step chain of errors (FM25): (1) `bar differential is closed color' WRONG --- $d_B$ is single degree-1 map over $(\ChirAss)^!$, not parametrized. (2) `deconcat = open color' WRONG --- cofree coassociative on $T^c(V)$; single-colored $E_1$. (3) `$d_B + \Delta = $ SC-coalgebra' WRONG --- two structures $\ne$ two colors. (4) `SC on $B(A)$ dual to SC on $(Z^{\mathrm{der}}(A), A)$' WRONG --- $B(A)$ is INPUT; SC emerges in OUTPUT via $\Hom(B(A), A)$. (5) `$E_1$-chiral $= E_1$-topological' WRONG --- chiral is on CURVE (2-real-dim holomorphic); topological is $\Conf_k(\R)$; bar is over $(\ChirAss)^!$, not $(\Ass)^!$. (6) Steinberg analogy RETIRED.''
\end{quote}

Reading this against Theorem~\ref*{thm:bar-diff-eq-holfact}: items (1) and (2) of the error chain are precisely what Theorem~\ref*{thm:bar-diff-eq-holfact} \emph{appears} to assert if read carelessly. The chapter avoids items (1) and (2) mathematically, by inserting the factorisation-algebra interpretation between the bar complex and the coloured operad; but the prose near l.\,388 could land in item (3) if not read with the full proof in hand.

\paragraph{Heal 6c (synthesis).} The mathematically honest articulation is:
\begin{enumerate}[label=(\arabic*)]
\item $B^{\mathrm{ch}}(\cA) = T^c(s^{-1}\bar\cA)$ is an $E_1$ coassociative coalgebra. Single structure, single colour.
\item The bar differential $d_B$ is a chain differential on $B^{\mathrm{ch}}(\cA)$, constructed from the chiral product; it identifies under BD/CG with the holomorphic factorisation map $\mu^{\mathrm{hol}}_\Delta$ on the factorisation algebra $\Obs^{\mathrm{ch}}_\cA$.
\item The coproduct $\Delta_B$ is the cofree deconcatenation on $T^c(s^{-1}\bar\cA)$; it identifies under BD/CG with the topological factorisation map $\mu^{\mathrm{top}}$ on $\Obs^{\mathrm{ch}}_\cA|_{\Ran(\R)}$.
\item The pair $(Z^{\mathrm{der}}_{\mathrm{ch}}(\cA), \cA) \simeq (\Obs^{\mathrm{ch}}_\cA|_{\Ran(X)}, \Obs^{\mathrm{ch}}_\cA|_{\Ran(\R)})$ carries the $\SCchtop$-coloured-operad action. The bar complex $B^{\mathrm{ch}}(\cA)$ is the Koszul-dual chain model of this pair; it \emph{presents} the pair via bar--cobar, but it is not itself an $\SCchtop$-coalgebra.
\item Theorem~\ref*{thm:bar-diff-eq-holfact} asserts: the bar differential and coproduct on $B^{\mathrm{ch}}(\cA)$ are the chain-level realisations of the two factorisation maps of the $\SCchtop$-datum $(Z^{\mathrm{der}}_{\mathrm{ch}}(\cA), \cA)$, and they satisfy the compatibility~\eqref*{eq:bar-coalg-mc} as a consequence of the factorisation-algebra structure.
\end{enumerate}

Under this articulation, AP165 is respected; no error of type (3) or (4) occurs; the prose of Theorem~\ref*{thm:bar-diff-eq-holfact} can be sharpened by making item (4) above explicit.

\textbf{Status of the heal.} AP165 correction is substantively propagated in \texttt{sc\_chtop\_heptagon.tex}; a minor prose residue in Theorem~\ref*{thm:bar-diff-eq-holfact}'s concluding clause is an improvement target, not a mathematical error. The chapter's contents respect AP165; the prose articulation can be sharpened.

%%%%%%%%%%%%%%%%%%%%%%%%%%%%%%%%%%%%%%%%%%%%%%%%%%%%%%%%%%%%%%%%%%%%%%%%%%%%%
\section{Verification on \texorpdfstring{$V_k(\fsl_2)$}{V\_k(sl\_2)} at \texorpdfstring{$k = 1$}{k=1}: the simplest nontrivial conformal witness}
%%%%%%%%%%%%%%%%%%%%%%%%%%%%%%%%%%%%%%%%%%%%%%%%%%%%%%%%%%%%%%%%%%%%%%%%%%%%%

Per the task prompt: ``Verify all seven on $V_k(\fsl_2)$ at $k = 1$.'' Level $k = 1$ is non-critical ($k + h^\vee = 3 \ne 0$) and minimal; central charge $c = 3k/(k + 2) = 1$; Sugawara tensor $T = \tfrac{1}{6}(e f + f e + \tfrac{1}{2} h h)$. Shadow class L; $\kappa = \dim(\fsl_2) \cdot (k + h^\vee)/(2 h^\vee) = 3 \cdot 3 / 4 = 9/4$ at Vol~I trace-form normalisation.

\paragraph{Face 1 (operadic).} $\FM_k(\mathbb H) \times_{\Conf_m(\R)} \Conf_m(\R)$ with $k$ closed insertions of $V_1(\fsl_2)$-vectors and $m$ open insertions of WZW-boundary modules. The component $\SCchtop(2\mathsf{cl}, 0\mathsf{op}; \mathsf{cl}) = \FM_2(\C)$ carries the $V_1(\fsl_2)$-OPE collision at coincident points, encoded as the Frenkel--Kac free-field realisation
\[
e(z) \leftrightarrow e^{i\alpha(z)}, \qquad f(z) \leftrightarrow e^{-i\alpha(z)}, \qquad h(z) \leftrightarrow i\partial\alpha(z)
\]
with $\alpha$ the $\mathrm{II}_{1, 1}$ lattice boson at self-dual radius. The collision stratum is 1-dimensional (single first-order pole), matching $\dim\FM_2(\C) = 1$ as an $S^1$-bundle.

\paragraph{Face 2 (Koszul dual).} $(\SCchtop)^! = (\Lie, \Ass, \text{shuffle-mixed})$ evaluated on $V_1(\fsl_2)$: closed sector is $\fsl_2$ as a Lie algebra (with Killing form, non-degenerate); open sector is the associative algebra of WZW boundary modules at level 1 (Chern--Simons line-bundle fusion ring, $\mathrm{SU}(2)_1 = \{1, \sigma\}$ with $\sigma \otimes \sigma = 1$); shuffle-mixed is the cross-term combining affine fusion with associative composition. The free-field realisation reduces the closed-sector Lie bracket to the $\mathrm{II}_{1, 1}$ lattice commutator $[\alpha_m, \alpha_n] = m \delta_{m + n, 0}$.

\paragraph{Face 3 (factorisation).} $\Obs^{\mathrm{ch}}_{V_1(\fsl_2)}$ on $\Ran(X) \times \Ran_{\ge 0}(\R)$. The closed-colour holomorphic factorisation $\mu^{\mathrm{hol}}_\Delta$ on $e, f$ reads (via the standard $\hat{\fsl}_2$ OPE at level~1)
\[
e(z) f(w) \sim \frac{1}{(z - w)^2} + \frac{h(w)}{z - w}, \quad\Rightarrow\quad \mu^{\mathrm{hol}}_\Delta(e, f) = \Delta_*[(e_{(-1)} f)(w)] = \Delta_*[(ef)(w)],
\]
the chiral (normally ordered) product. The central coefficient $e_{(1)} f = k = 1$ at level 1 is the $1/(z - w)^2$ coefficient; the bracket $e_{(0)} f = h$ is the $1/(z - w)$ coefficient. The open-colour topological factorisation $\mu^{\mathrm{top}}$ on two boundary insertions $(a, b)$ on $\R$ at $t_1 < t_2$ reads
\[
\mu^{\mathrm{top}}(a \otimes b) = (s^{-1} a) \otimes (s^{-1} b),
\]
the deconcatenation along the half-line.

\paragraph{Face 4 (BV/BRST).} Chern--Simons BV at level 1 on $X \times \R_{\ge 0}$ with boundary $\fsl_2$-WZW at level 1. The action is free quadratic (Chern--Simons at level 1 has elliptic kernel $\partial_{\bar z} \partial_t$ plus level term); boundary condition is Dirichlet-compatible (WZW currents at $t = 0$); CG axioms verified (class L scope). The propagator $1/(k + h^\vee) = 1/3$ is finite. The Costello--Gaiotto holographic picture (arXiv:1810.01970) makes this explicit: the 3d HT theory is the holographic dual of the 2d WZW boundary theory.

\paragraph{Face 5 (convolution $L_\infty$).} $\fg^{\SCchtop}_{V_1(\fsl_2), Z} = \Hom^{\mathrm{col}}(B(\SCchtop), \End^{\mathrm{ch}}_{(V_1(\fsl_2), Z)})$. The weight-2 MC element is generated by the classical $r$-matrix
\[
r^{KM}(z) = k \cdot \Omega / z = \Omega / z \quad \text{at } k = 1,
\]
with $\Omega = \sum_a T_a \otimes T^a$ the Casimir tensor of $\fsl_2$. The weight-3 MC closure is the $\fsl_2$-Jacobi identity (finite-dim classical $r$-matrix, no higher orders).

\paragraph{Face 6 (Drinfeld centre).} $\Rep_{\mathrm{fact}}(V_1(\fsl_2))$ is the level-1 affine Kazhdan--Lusztig category $\mathrm{KL}_1(\fsl_2)$, equivalent (after Finkelberg~1996) to $\mathrm{Rep}(U_q(\fsl_2))$ at $q = -1$ a primitive $(k + 2) = 3$rd root of unity. The two simple objects are the trivial module $V_0$ and the spin-$1/2$ module $V_1$; fusion is $V_1 \otimes V_1 = V_0$. The Drinfeld centre
\[
\cZ(\mathrm{KL}_1(\fsl_2)) \simeq \mathrm{KL}_1(\fsl_2) \boxtimes \mathrm{KL}_1(\fsl_2)^{\mathrm{rev}}
\]
by Müger's theorem (modular tensor categories are non-degenerate); this has 4 simple objects. The half-braiding is the Chern--Simons monodromy
\[
\sigma_{V_1(\fsl_2)}(z)|_{V_1 \otimes V_1} = R_{\mathrm{CS}} = \exp(i\pi \Omega / (k + 2)) = \exp(i\pi \Omega / 3),
\]
the fundamental $R$-matrix at $q = e^{i\pi/3}$. Spectral: $R(u - v) = $ the hyperbolic $\mathrm{SU}(2)_1$ $R$-matrix at rapidity $u - v$.

\paragraph{Face 7 (derived AG).} $\cM_{\SCchtop, X}$ at the point $(V_1(\fsl_2), Z^{\mathrm{der}}_{\mathrm{ch}}(V_1(\fsl_2)))$ is locally smooth with tangent dg-Lie of total $\Z$-graded dimension 5 (as recorded in Remark~\ref*{rem:heptagon-witness-vk-sl2}, l.\,2047): 3 classes from the $\fsl_2$-Casimir tower plus 2 from the chiral $r$-matrix correction. The 1-shifted symplectic form restricts to the Killing pairing on the $\fsl_2$-trace, non-degenerate (semisimple locus, Face~7 scope). The Lagrangian MC section carries $r^{KM}(z) = \Omega/z$ at $k = 1$.

\paragraph{Closure.} Traversing the heptagon (1) $\to$ (2) $\to$ (3) $\to$ (4) $\to$ (5) $\to$ (6) $\to$ (7) $\to$ (1) on $V_1(\fsl_2)$ returns to Face~1 via the Drinfeld-associator-twist $\Phi \in \mathrm{GRT}_1(\Q)$ (the Drinfeld--Kohno associator at level 1 reducing to the universal $\mathrm{SU}(2)_1$ braid representation). The triple-check on the closure:
\begin{itemize}
\item \emph{Operadic dimension}: $\FM_2(\C) = S^1$ (dim 1); the full strata of $\FM_k(\mathbb H) \times_{\Conf_m(\R)} \Conf_m(\R)$ at $(k, m) = (2, 2)$ give the 5 expected generators by the $\fsl_2$ Verlinde formula at level 1.
\item \emph{Centre dimension}: $\cZ(\mathrm{KL}_1(\fsl_2))$ at the Casimir-generated locus has 5 simple objects: the 4 from $\mathrm{KL}_1 \boxtimes \mathrm{KL}_1^{\mathrm{rev}}$ plus one orbifold-sector identification from the $\mathrm{SL}_2(\Z)$ action.
\item \emph{Tangent dimension}: total $\Z$-graded dimension 5 (Remark~\ref*{rem:heptagon-witness-vk-sl2}).
\end{itemize}
Three independent counts $= 5$. The heptagon closes on $V_1(\fsl_2)$ at the simplest non-trivial conformal level.

\paragraph{Level-$k$ cross-check.} At $k = 1$ the central charge is $c = 1$; at $k = 2$ one has $c = 3/2$; at general $k \ne -2$, $c = 3k/(k + 2)$. The heptagon closure is $k$-independent: the tangent dimension 5 is a topological invariant (class L with trivial super-structure), the Drinfeld centre is a level-invariant categorical construction on modular tensor categories (Müger), and the closed-operadic-dimension count is geometric. The $k$-dependence enters only through the central charge $c$ at the level of Face~3's Sugawara $T(z)$, which modifies the 3-pole coefficient of $r^{\mathrm{Vir}}(z) = (c/2)/z^3 + 2T/z$, but does not modify the heptagon closure count itself.

%%%%%%%%%%%%%%%%%%%%%%%%%%%%%%%%%%%%%%%%%%%%%%%%%%%%%%%%%%%%%%%%%%%%%%%%%%%%%
\section{Synthesis and status}
%%%%%%%%%%%%%%%%%%%%%%%%%%%%%%%%%%%%%%%%%%%%%%%%%%%%%%%%%%%%%%%%%%%%%%%%%%%%%

\textbf{Summary of the six attack axes.}
\begin{enumerate}[label=(\alph*)]
\item Seven presentations with sources: clarified in cycle 1 (table). No mathematical issues; prose sharpening possible.
\item 5/7 chain-level, 2/7 cohomological: cycle 2 identified Faces~6 (Drinfeld centre) and 7 (DAG) as the cohomological-only faces, driven by intrinsic $(\infty, 2)$-categorical and derived-stack universal properties.
\item Brace action --- Kadeishvili or chiral refinement: cycle 3, \emph{chiral refinement} specialising to Kontsevich at formal disc (three enhancements: spectral parameter, curve-global descent, mixed-sector compatibility).
\item Drinfeld-centre as ninth face: cycle 4, no. Two separate $\mathrm{GRT}_1$-torsors are at play; the ``nine faces'' is the $r(z)$-torsor (Vol~II Part III) and the heptagon is seven faces of $\SCchtop$-as-operad.
\item Derived-AG face --- equivalence or weaker: cycle 5, category equivalence of derived stacks; pointwise + Kan-extended at the operadic level; three scope qualifiers (semisimple locus, directional restriction, no-conformal-vector stratum).
\item AP165 full propagation: cycle 6, substantively propagated. Mathematical content respects AP165; one minor prose residue at Theorem~\ref*{thm:bar-diff-eq-holfact}'s concluding clause is an improvement target, not a mathematical error.
\end{enumerate}

\textbf{V_k($\fsl_2$) at $k = 1$ closure witnessed.} Three independent dimension counts (operadic, categorical centre, tangent) all equal 5. The heptagon is chain-level-closed on this simplest non-trivial conformal witness after fixing the Drinfeld--Kohno associator $\Phi \in \mathrm{GRT}_1(\Q)$.

\textbf{One quality-improvement target.} Rephrase the concluding clause of Theorem~\ref*{thm:bar-diff-eq-holfact} (l.\,388--399) to make explicit that $B^{\mathrm{ch}}(\cA)$ is the Koszul-dual chain model of the $\SCchtop$-datum pair, not an $\SCchtop$-coalgebra in its own right. The mathematical content is unchanged; the prose becomes AP165-explicit rather than AP165-implicit.

\textbf{One conceptual clarification flagged.} The programme-level architectural docs (\texttt{platonic\_ideal\_reconstituted\_2026\_04\_17.md}) collapse two distinct torsors --- the $\SCchtop$-operadic heptagon and the $r(z)$-$\mathrm{GRT}_1$-torsor --- under the shared word ``seven''. A manuscript reader could profit from a single sentence at the start of the heptagon chapter that distinguishes the two.

\textbf{No open mathematical errors found.} All six attack axes heal with the chapter's content intact. The correction work of AP165 is substantive, not superficial.
